\writeSectionFrame{\presentationTitle}{Unveränderliche (immutable) Objekte}
\begin{frame}{Idee: Immutable Objects}
    \begin{block}{Definition}
        Ein immutable object ist ein Objekt, dessen Zustand sich nach
        der Erzeugung nicht mehr ändern kann.
    \end{block}
    \pause
    \begin{itemize}
        \item Zustand wird nur im Konstruktor gesetzt
        \item Keine Setter
        \item Änderungen erzeugen neue Objekte
    \end{itemize}
\end{frame}

\begin{frame}[fragile]{Immutable Object – Beispiel}
    \begin{exampleblock}{Monster.java}
        \begin{lstlisting}
public class Monster {
    private final int health;

    public Monster(int health) {
        this.health = health;
    }

    public Monster takeDamage(int damage) {
        return new Monster(health - damage);
    }

    public int getHealth() {
        return health;
    }
}
        \end{lstlisting}
    \end{exampleblock}
\end{frame}

\begin{frame}[fragile]{Immutable Object – Beispiel}
    \begin{exampleblock}{Player.java}
        \begin{lstlisting}[basicstyle=\tiny]
public class Player {
	private final int health;
	private final int gold;

    public Player(int health, int gold) {
        this.health = health;
        this.gold = gold;
    }
    
    public Player takeGold(int amount) {
    	return new Player(health, this.gold + amount);
    }
    
    public Player spendGold(int amount) {
    	return new Player(health, this.gold - amount);
    }

    public Player takeDamage(int damage) {
        return new Player(health - damage, gold);
    }

    public int getHealth() {
        return health;
    }
}
        \end{lstlisting}
    \end{exampleblock}
\end{frame}

\begin{frame}[fragile]{Wir haben ein Problem}
    \begin{alertblock}{Player.java}
        \begin{lstlisting}
public class Player {
	// ...

    public Player attack(Monster monster) {
    	// ????
    }
}        
        \end{lstlisting}
    \end{alertblock}
    \begin{alertblock}{Warum?}
    	\begin{itemize}
    	  \item Zwei Objekte ändern logisch ihren Zustand
    	  \item Immutable Objekte können ihren Zustand nicht ändern
    	  \item Eine Methode kann aber nur einen Rückgabewert haben
    	\end{itemize}
    \end{alertblock}
\end{frame}

\begin{frame}{Überlegungen zur Lösung}
	\begin{itemize}
	  \item Immutable Objects erzwingen saubere Trennung von
	  	\begin{itemize}
	  	  \item Berechnung
	  	  \item Zustandsübergang
	  	  \item Koordination
	  	\end{itemize}
	  \item Ein Immutable Object sollte nicht mehrere Aggregate verändern
	  \item Kein technisches Problem, sondern ein Modellierungs­problem
	\end{itemize}
\end{frame}

\begin{frame}{Lösungsmöglichkeiten}
	\begin{itemize}
	  \item Monster veränderlich definieren
	  \item Ergebnisobjekt zurückgeben
	  \item Die Aktion aus dem Entity herausziehen (DDD-Stil)
	\end{itemize}
\end{frame}

\begin{frame}{Monster veränderlich definieren}
	\begin{alertblock}{Keine Lösung}
		\begin{itemize}
		  \item bricht Immutability
		  \item versteckter Seiteneffekt
		  \item genau das, was wir vermeiden wollten
		\end{itemize}	
	\end{alertblock}
\end{frame}

\begin{frame}{Ergebnisobjekt zurückgeben}
	\begin{itemize}
	  \item Man führt ein Result- / Event- / Outcome-Objekt ein
	  \item Ergebnis wird klar von der Aktion getrennt
	  \item Objekte können unveränderlich bleiben
	\end{itemize}	
\end{frame}

\begin{frame}[fragile]{Ergebnisobjekt zurückgeben}
	\begin{exampleblock}{AttackResult.java}
		\begin{lstlisting}
public class AttackResult {
    private final Player player;
    private final Monster monster;

    public AttackResult(Player player, Monster monster) {
        this.player = player;
        this.monster = monster;
    }

    public Player getPlayer() {
        return player;
    }

    public Monster getMonster() {
        return monster;
    }
}		
		\end{lstlisting}
	\end{exampleblock}
\end{frame}

\begin{frame}[fragile]{Neue attack-Methode}
	\begin{exampleblock}{Player.java}
		\begin{lstlisting}
public AttackResult attack(Monster monster) {
    Monster damagedMonster = monster.takeDamage(
    	this.attackPower);
    Player damagedPlayer = this.takeDamage(
    	monster.getAttackPower());

    return new AttackResult(damagedPlayer, damagedMonster);
}		
		\end{lstlisting}
	\end{exampleblock}
\end{frame}

\begin{frame}{Bewertung}
	\begin{itemize}
	  \item Beide neuen Zustände sind explizit sichtbar
	  \item Kein Objekt wird verändert
	  \item Sehr gut testbar
	  \item Sehr modernes Pattern
	\end{itemize}	
\end{frame}

\begin{frame}{Die Aktion aus dem Entity herausziehen}
	\begin{alertblock}{Frage}
		Gehört die attack-Methode überhaupt in ein Entity?	
	\end{alertblock}
	\pause
	
	\begin{exampleblock}{Antwort}
		Man könnte sie auch als Bestandteil der Geschäftslogik betrachten!
	\end{exampleblock}
\end{frame}

\begin{frame}[fragile]{Klasse der Geschäftslogik}
	\begin{exampleblock}{CombatServiceImpl.java}
		\begin{lstlisting}
public class CombatServiceImpl implements CombatService {
	@Override
    public AttackResult attack(Player player, 
    		Monster monster) {
        Monster newMonster = monster.takeDamage(
        	player.getAttackPower());
        Player newPlayer = player.takeDamage(
        	monster.getAttackPower());

        return new AttackResult(newPlayer, newMonster);
    }
}
		\end{lstlisting}
	\end{exampleblock}
\end{frame}

\begin{frame}{Bewertung}
	\begin{itemize}
	  \item Die Entity-Klassen bleiben schlank und enthalten nur Daten und Regeln
	  \item Kampf ist eine Domänenaktion, kein Objektproblem
	  \item Das ist Clean Architecture pur (Domain Driven Design)
	  \item Sehr moderne Softwarearchitektur
	\end{itemize}	
\end{frame}